\documentclass[11pt]{amsart}
\usepackage{amsmath,amssymb,amsthm,mathtools}
\usepackage[margin=1in]{geometry}
\usepackage{hyperref}
\usepackage{listings}
\lstset{basicstyle=\ttfamily\footnotesize,breaklines=true,frame=single,
  showstringspaces=false,columns=fullflexible,keepspaces=true}

\newtheorem{theorem}{Theorem}[section]
\newtheorem{proposition}[theorem]{Proposition}
\newtheorem{lemma}[theorem]{Lemma}
\newtheorem{corollary}[theorem]{Corollary}
\theoremstyle{definition}
\newtheorem{definition}[theorem]{Definition}
\theoremstyle{remark}
\newtheorem{remark}[theorem]{Remark}
\newtheorem{example}[theorem]{Example}

\newcommand{\RH}{\mathrm{RH}}
\newcommand{\TP}{\mathrm{TP}}

\title[A Uniformity Gap]{Finite-to-full positivity is a uniformity gap: a criterion, a
Hodge--Riemann stability theorem for towers, and the place of $\zeta$}
\author{David Alejandro Trejo Pizzo}
\thanks{Independent Researcher. \texttt{dtrejopizzo@gmail.com}}


\begin{document}
\nocite{bgIK,bgMV,bgTitchmarsh,bgConrey03}
\begin{abstract}
For the Riemann $\xi$-function every finite-order positivity --- the Tur\'an
inequalities, the hyperbolicity of each fixed-degree Jensen polynomial, each
finite total-positivity minor --- is unconditional, while the full-order
positivity is the Riemann Hypothesis. We isolate the general phenomenon as a
\emph{uniformity} (quantifier-order) gap, not a logical-implication gap, and
develop the resulting positivity theory. We prove a criterion: the bridge from
``eventually at each order'' to ``uniformly at all orders'' holds precisely when
the object is scale-free or has a bounded order-threshold; we show the
bounded-threshold class is non-empty (converging strictly-positive measures) and
that its discriminator is uniform convergence to a definite limit; we prove that
the convergence route to RH is over-shooting or circular, so the framework
explains the gap without crossing it; and we prove a \emph{Hodge--Riemann
stability theorem}: a Hodge--Riemann tower has uniform positivity if and only if
its limit form is strictly definite, generalizing Adiprasito--Huh--Katz from a
single object to families. The verdicts place $\zeta$'s hierarchy at the
borderline degenerate limit. A numerical appendix with code illustrates the
stability dichotomy. Everything is independent of the Riemann Hypothesis.
\end{abstract}
\maketitle

\tableofcontents
\newpage

\section{Introduction}
Several reformulations of RH are full-order positivity statements whose
finite-order truncations are unconditional: the total positivity of the de
Bruijn--Newman kernel, the hyperbolicity of the Jensen polynomials of $\xi$, the
Hankel (moment) positivity. Across these, every finite order is known
unconditionally --- for each fixed order, eventually in height --- while the full
order is RH~\cite{RP13,RP15}. This paper extracts the structure responsible,
shows it is a uniformity gap with a clean criterion, and develops the attendant
positivity theory, culminating in a stability theorem for Hodge--Riemann towers.
The results are independent of RH; the connection to $\zeta$ is that its
hierarchy sits at the borderline of the criterion.

\section{Scale-parametrized positivity hierarchies}
\begin{definition}
A \emph{scale-parametrized positivity hierarchy} is a family
$\{P_k(T)\}_{k\ge1,\,T\ge1}$ of conditions, nested $P_1\supseteq P_2\supseteq
\cdots$, with a scale parameter $T$, together with order defects
$\delta_k(T)\ge0$ such that $P_k(T)\iff\delta_k(T)\le0$.
\end{definition}
Two statements must be separated:
\[
\textbf{(eventual)}\quad\forall k\ \exists T_0(k)\ \forall T>T_0(k):P_k(T);
\qquad
\textbf{(uniform)}\quad\forall T\ \forall k:P_k(T).
\]
The full order $P_\infty=\bigcap_k P_k$ holding at every scale is \textbf{(uniform)};
the known unconditional results are \textbf{(eventual)}.

\begin{remark}
The finite-to-full gap is \emph{not} the implication $\forall k\,P_k\Rightarrow
P_\infty$, which is trivial since $P_\infty=\bigcap_kP_k$. It is the failure of
the quantifier swap: each $P_k$ is known only eventually in $T$, with threshold
$T_0(k)\to\infty$, so no single scale enjoys all orders. For $\zeta$ the orders
are: $\TP_1$ ($\Phi>0$), $\TP_2$ (Tur\'an/log-concavity/degree-two Jensen),
order $k$ (degree-$k$ Jensen hyperbolicity, Griffin--Ono--Rolen--Zagier), and
$\TP_\infty=$ RH.
\end{remark}

\section{The uniformity criterion}
\begin{theorem}[Criterion]\label{thm:crit}
For a scale-parametrized hierarchy, \textbf{(eventual)} implies
\textbf{(uniform)} if and only if $\sup_k\sup_T\delta_k(T)\le0$; this holds in
either of two cases:
\begin{enumerate}
\item \emph{finite type:} the object is scale-free, $\delta_k(T)=\delta_k$
independent of $T$ (so eventual $=$ uniform); or
\item \emph{bounded threshold:} $\sup_kT_0(k)<\infty$.
\end{enumerate}
\end{theorem}
\begin{proof}
\textbf{(uniform)} is exactly $\sup_k\sup_T\delta_k(T)\le0$. In case (1) the
inner supremum over $T$ is $\delta_k$, and \textbf{(eventual)} forces
$\delta_k\le0$. In case (2), for $T>\sup_kT_0(k)$ every $P_k(T)$ holds; the
finitely many smaller scales are covered by \textbf{(eventual)} at each fixed
$k$. Conversely, if both fail, some $\delta_k(T)>0$ persists for arbitrarily
large $T$, contradicting \textbf{(uniform)}.
\end{proof}

\section{The dividing line}
\begin{center}
\begin{tabular}{lccc}
\hline
hierarchy & scale parameter? & threshold $T_0(k)$ & finite-to-full \\
\hline
matroid Hodge--Riemann (AHK) & none & --- & \textbf{bridge holds} \\
Lorentzian / stable polynomials & none & --- & \textbf{bridge holds} \\
Hamburger moment problem & none & --- & bridge holds \\
$\zeta$ Jensen / total positivity & height $T$ & $T_0(k)\to\infty$ & \textbf{gap} \\
\hline
\end{tabular}
\end{center}
Finite combinatorial or fixed-measure objects are scale-free (case 1), so their
full positivity follows from a general theorem (Adiprasito--Huh--Katz for
matroids). The $\zeta$-hierarchies carry a genuine height parameter with
$T_0(k)\to\infty$, so neither case applies.

\section{Equivalence with the finitization obstruction}
\begin{corollary}\label{cor:fin}
A scale-parametrized hierarchy has no finite-to-full gap iff it is finite type or
bounded threshold. For the $\zeta$-hierarchies, a finite-dimensional model (a
Frobenius-type periodicity of the zeros) would make the object scale-free, hence
case (1); its absence~\cite{RP15} is the absence of case (1), and the unbounded
thresholds are the failure of case (2). Thus
\[
\text{finite-to-full gap}\ \Longleftrightarrow\ \text{neither finite type nor bounded threshold}\ \Longleftrightarrow\ \text{no finitization}.
\]
\end{corollary}

\section{The bounded-threshold class is non-empty}
\begin{proposition}[Existence]\label{prop:exist}
There exist scale-parametrized positivity hierarchies that are not finite type
yet have bounded threshold.
\end{proposition}
\begin{proof}
Let $\mu_T\to\mu_\infty$ in total variation, $\|\mu_T-\mu_\infty\|\le
\varepsilon(T)\to0$, with $\mu_\infty$ strictly positive (every moment/total-
positivity minor $>0$, with a uniform gap $g_0>0$ from the boundary on the unit
sphere of test vectors). Each finite minor is a continuous functional of the
measure, $1$-Lipschitz in total variation after normalization; hence for
$T>T^\ast:=\sup\{T:\varepsilon(T)>g_0\}$ every minor of $\mu_T$ is positive, so
$\sup_kT_0(k)\le T^\ast<\infty$, while $\mu_T$ genuinely varies with $T$.
\end{proof}
The discriminating mechanism is therefore \emph{uniform convergence to a strictly
positive limit}. The $\xi$-Jensen polynomials converge (rescaled) to Hermite
polynomials --- a hyperbolic limit --- but the Griffin--Ono--Rolen--Zagier
thresholds $T_0(k)$ grow: the convergence is not uniform in degree, so $\zeta$
sits outside the bounded-threshold class.

\section{The convergence route is over-RH or circular}
One might hope the convergence mechanism (not obstructed by linear independence,
unlike a finitization) provides a passage to RH. It does not.
\begin{theorem}\label{thm:overRH}
Uniform $(d,n)$ convergence of the Jensen polynomials to the rescaled Hermite
limit at rate $o(1/d)$ implies RH (the Hermite gap is $\asymp1/d$, so being
within that margin gives hyperbolicity for all $(d,n)$); but it is not implied by
RH (a hyperbolic $J_{d,n}$ need not match the Hermite spacing). Hence uniform
convergence is strictly stronger than RH. The RH-equivalent relaxation ---
``$J_{d,n}$ within the margin of \emph{some} hyperbolic polynomial'' --- is just
``$J_{d,n}$ hyperbolic,'' i.e.\ RH itself.
\end{theorem}
\begin{proof}
The rescaled Hermite limit $\widehat H_d$ has simple real zeros with minimal gap
$\asymp1/d$; a real perturbation of size $o(1/d)$ in coefficients moves each zero
by $o(1/d)$ and preserves real-rootedness, so $J_{d,n}$ is hyperbolic; over all
$(d,n)$ this is RH. Conversely RH asserts hyperbolicity, a property of the zero
set's reality, not of its proximity to the Hermite configuration; a hyperbolic
polynomial may have any real zero spacing. The relaxation to ``some hyperbolic
polynomial'' is hyperbolicity by definition.
\end{proof}
So the convergence route over-shoots RH (harder) or restates it (circular); there
is no intermediate sub-RH opening. The framework's RH-relevance is exactly to
explain the gap, not to cross it.

\section{A Hodge--Riemann stability theorem for towers}
We now develop the positivity theory in its own right, generalizing
Adiprasito--Huh--Katz from a single object to a family.

\begin{definition}[Hodge--Riemann tower]
A \emph{Hodge--Riemann tower} is a sequence $(V_m,L_m,Q_m)$, where $V_m=
\bigoplus_jV_m^j$ is a finite graded real vector space, $L_m\colon V_m^j\to
V_m^{j+1}$ a Lefschetz operator, and $Q_m$ the bilinear form on primitive parts
$P_m^j=\ker L_m^{\,d-2j+1}$, with the Hodge--Riemann property
$(-1)^jQ_m|_{P_m^j}>0$. Set the \emph{HR gap}
$\delta_m=\inf_{j}\inf_{\|v\|=1,\,v\in P_m^j}(-1)^jQ_m(v,v)$. The tower
\emph{converges} to $(V_\infty,L_\infty,Q_\infty)$ if, on a common model (fixed
grading and dimensions, or an inverse limit), $Q_m\to Q_\infty$ and the primitive
projections converge in operator norm.
\end{definition}

\begin{theorem}[HR stability]\label{thm:hr}
Let $(V_m,L_m,Q_m)$ be a convergent Hodge--Riemann tower with limit gap
$\delta_\infty$.
\begin{enumerate}
\item If $\delta_\infty>0$ (the limit form is strictly definite on primitives),
the tower has \emph{uniform Hodge--Riemann}: $\liminf_m\delta_m\ge\delta_\infty>0$,
so $\delta_m\ge\delta_\infty/2$ for all large $m$.
\item If $\delta_\infty=0$ (the limit form degenerates on some primitive
direction), uniform HR fails: $\inf_m\delta_m=0$.
\end{enumerate}
\end{theorem}
\begin{proof}
On the common model, $m\mapsto Q_m$ and $m\mapsto P_m^j$ (as orthogonal
projections) are convergent, so $(m,v)\mapsto(-1)^jQ_m(v,v)$ is jointly continuous
on the compact unit sphere of $\bigoplus_jP_m^j$; hence $\delta_m\to\delta_\infty$.
If $\delta_\infty>0$, definiteness is an open condition: there is $m_0$ with
$\delta_m\ge\delta_\infty/2$ for $m\ge m_0$, giving (1). If $\delta_\infty=0$, a
unit primitive $v_\infty$ with $Q_\infty(v_\infty,v_\infty)=0$ is a limit of
$v_m\in P_m^{j}$ with $(-1)^jQ_m(v_m,v_m)\to0$, so $\inf_m\delta_m=0$, giving (2).
\end{proof}

\begin{corollary}\label{cor:towerbridge}
A Hodge--Riemann tower has the finite-to-full bridge (uniform positivity across
the tower) iff its limit form is strictly definite. This is the Hodge-theoretic
refinement of Theorem~\ref{thm:crit}(2): bounded threshold $=$ strictly-definite
limit.
\end{corollary}

\begin{remark}
Adiprasito--Huh--Katz prove HR for each individual matroid; Theorem~\ref{thm:hr}
governs \emph{families}, identifying strict-definiteness of the limit as the exact
condition for uniformity --- genuinely new content, since HR is classically a
per-object theorem. It predicts the dividing line of \S4: a hierarchy bridges
finite-to-full iff its objects converge to a strictly-definite limit. For
$\zeta$, the $\xi$-Jensen polynomials converge to Hermite (hyperbolic, so the
limit is definite in each fixed degree), but the convergence is non-uniform in
the degree direction, so the relevant limit gap across degrees degenerates: the
borderline $\delta_\infty=0$ case.
\end{remark}

\section{Worked examples}
\begin{example}[A uniform-HR tower]\label{ex:def}
Let $\mu_m$ be the probability measure $\tfrac1m\sum_{i=1}^m\delta_{i/m}$ on
$[0,1]$, converging weakly to Lebesgue measure $\mu_\infty$. The Hankel forms
$H_m^{(N)}=[\,\int x^{a+b}d\mu_m\,]_{0\le a,b\le N}$ converge to the Hilbert-type
Hankel matrix of $\mu_\infty$, which is strictly positive definite for every $N$
(Lebesgue measure has infinite support). By Proposition~\ref{prop:exist} the
total-positivity hierarchy of $\mu_m$ has bounded threshold: there is a single
$m_0(N)$, in fact $m_0$ \emph{uniform in $N$} on compact moment ranges, beyond
which all minors are positive. This is a bounded-threshold, non-finite-type
hierarchy.
\end{example}

\begin{example}[The degenerate borderline]\label{ex:deg}
Let $\nu_m=(1-\tfrac1m)\delta_0+\tfrac1m\,\lambda_{[0,1]}$ converge to
$\nu_\infty=\delta_0$, a point mass. The limit Hankel form is rank one
(degenerate, $\delta_\infty=0$): all higher minors vanish. By
Theorem~\ref{thm:hr}(2) uniform positivity fails --- the order-$N$ minor is
$\asymp m^{-N}$, so $T_0(N)\to\infty$. This is the structural shape of $\zeta$'s
situation: a limit that is positive but \emph{degenerate} in the high-order
directions, so the thresholds grow.
\end{example}

\begin{example}[Uniform matroids]\label{ex:matroid}
The uniform matroids $U_{r,m}$ have characteristic polynomials whose
(normalized) coefficient sequences converge as $m\to\infty$; AHK gives HR for
each $U_{r,m}$ individually, and Theorem~\ref{thm:hr} adds that the family has
uniform HR iff the limiting normalized intersection form is definite --- a
concrete instance of the tower theorem in the combinatorial-Hodge setting where
AHK applies per object.
\end{example}

\section{The $\zeta$ orders, concretely}
For $\xi$ the hierarchy is the total positivity of the de Bruijn--Newman kernel
$\Phi$, with $\Phi>0$ ($\TP_1$) classical and the order-two defect
\[
\delta_2 = -\bigl(b(k)^2-b(k-1)b(k+1)\bigr),\qquad b(k)=m_{2k}/(2k)!,
\]
nonpositive (the Tur\'an inequality, hence $\TP_2$) unconditionally by
Csordas--Norfolk--Varga. The order-$k$ defect is the discriminant condition for
the degree-$k$ Jensen polynomial $J_{k,n}$, with threshold $T_0(k)=N(k)$ the
Griffin--Ono--Rolen--Zagier height beyond which $J_{k,n}$ is hyperbolic. The
content of \S6 is that $N(k)\to\infty$: the rescaled $J_{k,n}$ converge to the
Hermite polynomial $H_k$ (hyperbolic), but the convergence is non-uniform in $k$
because the Hermite gap shrinks like $1/k$ while the convergence rate at degree
$k$ does not keep pace --- the degenerate-borderline shape of
Example~\ref{ex:deg}, in the degree direction.

\section{Defect-rate transfer}
The criterion can be made quantitative.
\begin{proposition}[Rate transfer]\label{prop:rate}
Suppose the tower converges with $\|Q_m-Q_\infty\|\le\eta(m)$ and the limit is
strictly definite with gap $\delta_\infty>0$. Then $\delta_m\ge\delta_\infty-
C\eta(m)$ for a constant $C$ depending only on the primitive-projection moduli;
in particular uniform HR holds as soon as $\eta(m)<\delta_\infty/C$, with an
explicit threshold $m_0$. Conversely, if $\delta_\infty=0$ with
$\delta_m\asymp\eta(m)\to0$, the threshold $T_0(k)$ (in the scale language) grows
like the inverse rate, $T_0(k)\asymp\eta^{-1}(\text{gap at order }k)$.
\end{proposition}
\begin{proof}
$|\,(-1)^jQ_m(v,v)-(-1)^jQ_\infty(v,v)\,|\le\|Q_m-Q_\infty\|+2\|P_m^j-P_\infty^j\|
\,\|Q_\infty\|\le C\eta(m)$ uniformly over unit primitives; take infima to get
$|\delta_m-\delta_\infty|\le C\eta(m)$.
\end{proof}
This identifies the non-uniformity of the $\xi$-Jensen$\to$Hermite convergence as
the exact obstruction: the per-degree gap shrinks like $1/d$ while the
convergence rate at degree $d$ does not keep pace, so $T_0(d)\to\infty$.

\section{Conclusion}
The finite-to-full gap is a uniformity gap (Theorem~\ref{thm:crit}); its two
escape routes are finite type and bounded threshold, the latter a non-empty class
discriminated by uniform convergence to a definite limit
(Proposition~\ref{prop:exist}); the convergence route to RH over-shoots or
circles (Theorem~\ref{thm:overRH}); and Hodge--Riemann positivity is stable under
limits exactly when the limit form is definite (Theorem~\ref{thm:hr}). The
framework explains why $\zeta$'s hierarchy has the gap --- no finitization, a
degenerate limit gap across degrees --- without offering a passage to RH, and
stands on its own as positivity theory. The natural continuations, all RH-
independent, are the classification of uniform-HR towers, the construction of
bounded-threshold non-finite-type hierarchies beyond converging measures, and a
uniform Hodge--Riemann theorem for towers of matroids or polymatroids.

\appendix
\section{Numerical illustration of the stability dichotomy}\label{app:code}
The following program builds two convergent families of symmetric bilinear forms
on a fixed space --- one converging to a strictly definite limit, one to a
degenerate limit --- and tracks the gap $\delta_m$, illustrating
Theorem~\ref{thm:hr}.

\lstinputlisting[language=Python,caption={\texttt{hr\_stability\_demo.py}}]{code.py}

\medskip\noindent\textbf{Output.}
\begin{lstlisting}
 definite limit (delta_inf>0):  expect delta_m -> 0.5 (uniform HR holds)
   m=   2: delta=0.2010   m=   8: delta=0.4120   m=  32: delta=0.4767
   m= 128: delta=0.4941   m= 512: delta=0.4985
 degenerate limit (delta_inf=0):  expect delta_m -> 0 (uniform HR fails)
   m=   2: delta=0.5000   m=   8: delta=0.1250   m=  32: delta=0.0312
   m= 128: delta=0.0078   m= 512: delta=0.0020
\end{lstlisting}
The definite-limit tower has $\delta_m\to\delta_\infty=1/2>0$ (uniform HR); the
degenerate-limit tower has $\delta_m=1/m\to0$ (uniform HR fails), exactly as
Theorem~\ref{thm:hr} predicts, with the $1/m$ decay matching the rate transfer of
Proposition~\ref{prop:rate}.

\section{The degenerate borderline and threshold growth}\label{app:deg}
The degenerate case (Example~\ref{ex:deg}) is the structural shape of $\zeta$'s
hierarchy: a positive but \emph{degenerate} limit, so that higher-order
positivity holds only at a growing threshold. The following computes the
order-$N$ Hankel minimal eigenvalue of $\nu_m=(1-\tfrac1m)\delta_0+\tfrac1m
\lambda_{[0,1]}\to\delta_0$.
\lstinputlisting[language=Python,caption={\texttt{code2.py}: degenerate borderline}]{code2.py}
\noindent\textbf{Output.}
\begin{lstlisting}
   order-N Hankel smallest eigenvalue of nu_m (limit delta_0):
        m    N=2       N=3       N=4
        4  1.86e-03  5.19e-05  1.51e-06
       16  4.89e-04  1.32e-05  3.80e-07
       64  1.24e-04  3.32e-06  9.53e-08
      256  3.10e-05  8.31e-07  2.38e-08
     1024  7.74e-06  2.08e-07  5.96e-09
\end{lstlisting}
At fixed order the eigenvalue $\to0$ as $m\to\infty$ (the degenerate limit); at
fixed $m$ it drops by roughly three orders of magnitude from $N=2$ to $N=4$ ---
higher-order positivity is harder, so the threshold $T_0(N)$ grows in $N$. This is
the order-by-order signature of a degenerate limit, and the structural analogue,
in the moment hierarchy, of the Griffin--Ono--Rolen--Zagier threshold growth for
the $\xi$-Jensen polynomials.

\begin{thebibliography}{9}
\bibitem{RP13} D.~A.~Trejo Pizzo, \emph{Two modern routes and the wrong-sign obstruction}, preprint, 2026.
\bibitem{RP15} D.~A.~Trejo Pizzo, \emph{Why the Riemann Hypothesis resists}, preprint, 2026.
\bibitem{AHK} K.~Adiprasito, J.~Huh, E.~Katz, \emph{Hodge theory for combinatorial geometries}, Ann. of Math. \textbf{188} (2018).
\bibitem{GORZ} M.~Griffin, K.~Ono, L.~Rolen, D.~Zagier, \emph{Jensen polynomials for the Riemann zeta function and other sequences}, PNAS \textbf{116} (2019).
\bibitem{bgIK} H.~Iwaniec and E.~Kowalski, \emph{Analytic Number Theory}, Amer.\ Math.\ Soc.\ Colloq.\ Publ.\ \textbf{53}, AMS, 2004.
\bibitem{bgMV} H.~L.~Montgomery and R.~C.~Vaughan, \emph{Multiplicative Number Theory I: Classical Theory}, Cambridge Univ.\ Press, 2007.
\bibitem{bgTitchmarsh} E.~C.~Titchmarsh, \emph{The Theory of the Riemann Zeta-Function}, 2nd ed., rev.\ D.~R.~Heath-Brown, Oxford Univ.\ Press, 1986.
\bibitem{bgConrey03} J.~B.~Conrey, \emph{The Riemann Hypothesis}, Notices Amer.\ Math.\ Soc.\ \textbf{50} (2003), 341--353.
\end{thebibliography}

\end{document}
